\documentclass[11pt]{article}

\usepackage[a4paper,margin=2.5cm]{geometry}
\usepackage{amsmath,amssymb,amsthm,mathtools}
\usepackage{enumitem}
\usepackage{hyperref}

\title{Connected Radial Rigidity: What Survives from the Section 4 Strategy}
\author{}
\date{}

\newtheorem{theorem}{Theorem}[section]
\newtheorem{proposition}[theorem]{Proposition}
\newtheorem{lemma}[theorem]{Lemma}
\newtheorem{remark}[theorem]{Remark}

\newcommand{\C}{\mathbb C}
\newcommand{\R}{\mathbb R}
\newcommand{\1}{\mathbf 1}
\newcommand{\pb}{\partial_{\bar z}}

\begin{document}

\maketitle

\section{Purpose}

The simply connected argument in Section~4 of \texttt{main.tex} uses the holomorphic
auxiliary function
\[
W(z)=\pi z\left(4\partial_z u(z)-\frac{\Phi(|z|^2)}{z}+\frac{c}{\pi z}\right),
\]
whose boundary values are real on all of $\partial U$. In the connected but not simply
connected case, that argument no longer closes: reality on the \emph{outer} boundary does not
control the values of $\Im W$ on the inner boundary components.

What \emph{does} survive is the following scheme:
\begin{enumerate}[label=\textnormal{(\arabic*)}]
\item prove separately that the outer boundary component is a circle centered at the origin;
\item use the resulting collar to show that the angular derivative of the potential vanishes in
all of $U\setminus\{0\}$;
\item conclude that $U$ is rotationally invariant, hence a disk or an annulus.
\end{enumerate}

The first step is the genuinely hard one. The second and third steps are rigorous under mild
regularity assumptions and do not use simple connectivity.

\section{Outer Circle Implies Disk or Annulus}

Let $U\subset \C$ be bounded and connected. Denote by $\Gamma_{\mathrm{out}}$ its outer
boundary component, and by $E_\infty$ the unbounded connected component of
$\C\setminus \overline U$.

\begin{theorem}\label{thm:outer-circle-implies-annulus}
Assume $U\subset \C$ is bounded, connected, and has $C^1$ boundary. Let
$u\in C^1(\C\setminus\{0\})$ be compactly supported and satisfy
\begin{equation}\label{eq:pde}
\Delta u=\rho(|z|)\1_U-c\,\delta_0
\qquad\text{in }\mathcal D'(\C),
\end{equation}
where $\rho:[0,\infty)\to \R$ is continuous and
\[
\rho(r)>0 \qquad (r>0).
\]
Assume moreover that
\[
u\equiv 0 \qquad\text{in }E_\infty.
\]
If the outer boundary component is a circle centered at the origin,
\[
\Gamma_{\mathrm{out}}=\{|z|=R\},
\]
then $U$ is either a disk or an annulus centered at the origin.
\end{theorem}

\begin{proof}
Since the inner boundary components are compact and disjoint from $\Gamma_{\mathrm{out}}$, there
exists $\varepsilon>0$ such that the collar
\[
\mathcal A:=\{R-\varepsilon<|z|<R\}
\]
is contained in $U$.

Define the angular derivative
\[
L:=x\partial_y-y\partial_x=\partial_\theta
\]
and set
\[
w:=Lu.
\]
Because $L$ commutes with $\Delta$ and annihilates every radial function and the point mass
$\delta_0$, the equation \eqref{eq:pde} gives
\[
\Delta w=0
\qquad\text{in }U\setminus\{0\}.
\]
Thus $w$ is harmonic in $U\setminus\{0\}$.

On the outer boundary, the exterior side is identically zero and $u\in C^1(\C\setminus\{0\})$,
so
\[
u=0,\qquad \nabla u=0
\qquad\text{on }\Gamma_{\mathrm{out}}.
\]
Hence
\[
w=\partial_\theta u=0
\qquad\text{on }|z|=R,
\]
and also
\[
\partial_r w=\partial_r\partial_\theta u=\partial_\theta\partial_r u=0
\qquad\text{on }|z|=R.
\]

Since $w$ is harmonic in the annulus $\mathcal A$, it has the Fourier expansion
\[
w(r,\theta)=a_0+b_0\log r+\sum_{n\ge 1}(a_n r^n+b_n r^{-n})\cos(n\theta)
+\sum_{n\ge 1}(c_n r^n+d_n r^{-n})\sin(n\theta).
\]
The conditions
\[
w(R,\theta)=0,\qquad \partial_r w(R,\theta)=0
\]
force all coefficients to vanish, so
\[
w\equiv 0
\qquad\text{in }\mathcal A.
\]

Now $U\setminus\{0\}$ is connected: if $0\notin U$ there is nothing to prove, and if $0\in U$
then removing a single point from a planar domain leaves it connected. By unique continuation for
harmonic functions,
\[
w\equiv 0
\qquad\text{in }U\setminus\{0\}.
\]
Therefore
\[
\partial_\theta u=0
\qquad\text{in }U\setminus\{0\},
\]
so $u$ is radial there.

Applying $\Delta$ to a radial function gives a radial distribution. Since $\delta_0$ is radial,
\eqref{eq:pde} shows that
\[
\rho(|z|)\1_U
\]
is radial as a distribution. Because $\rho(r)>0$ for every $r>0$, the characteristic function
$\1_U$ must be radial on each circle $\{|z|=r\}$, hence $U$ is rotationally invariant about the
origin.

A bounded connected rotationally invariant planar domain is either a disk or an annulus centered
at the origin.
\end{proof}

\begin{remark}
This is the part of the connected-case strategy that is genuinely robust. No simple connectivity
is used anywhere in the proof.
\end{remark}

\section{A Conditional Hermite-Type Extension}

For the Hermite weights
\[
\rho_k(r):=r^{2k}e^{-\pi r^2},
\]
the Bargmann reduction gives the moment identity
\begin{equation}\label{eq:exp-identity}
\int_U e^{\pi zw}|z|^{2k}e^{-\pi |z|^2}\,dA(z)=c_k
\qquad (w\in \C).
\end{equation}
Expanding in powers of $w$ yields
\[
\int_U z^n |z|^{2k}e^{-\pi |z|^2}\,dA(z)=0
\qquad (n\ge 1).
\]

Define
\[
H_k(s):=s^{-(k+1)}\int_0^s t^k e^{-\pi t}\,dt.
\]
Then
\[
\pb\!\left(z^{n+k}\bar z^{\,k+1}H_k(|z|^2)\right)=z^n|z|^{2k}e^{-\pi|z|^2},
\]
so Cauchy--Green gives the boundary identity
\begin{equation}\label{eq:boundary-identity}
\int_{\partial U} z^{n+k}\bar z^{\,k+1}H_k(|z|^2)\,dz=0
\qquad (n\ge 1).
\end{equation}

\begin{theorem}\label{thm:conditional}
Assume the hypotheses of Theorem~\ref{thm:outer-circle-implies-annulus} with
\[
\rho(r)=r^{2k}e^{-\pi r^2},
\]
and assume in addition:
\begin{enumerate}[label=\textnormal{(\roman*)}]
\item the Bargmann moment identity \eqref{eq:exp-identity} holds;
\item the outer boundary component $\Gamma_{\mathrm{out}}$ is a real-analytic Jordan curve.
\end{enumerate}
Then $U$ is a disk or an annulus centered at the origin.
\end{theorem}

\begin{proof}
Under the analyticity assumption, the steepest-descent argument from the outer-boundary part of
\texttt{newnew.tex} becomes valid for $\Gamma_{\mathrm{out}}$: the contribution of the outer
boundary to \eqref{eq:boundary-identity} has asymptotic size
\[
R_*^n n^{-1/(2m_*)}
\left(\sum_{\ell=1}^L c_\ell e^{in\theta_\ell}\right)
\]
if $|\gamma(t)|$ is nonconstant on the outer boundary, whereas each inner boundary contribution is
only $O((R_j^*)^n)$ with $R_j^*<R_*$. Since a nontrivial finite exponential sum cannot tend to
zero, \eqref{eq:boundary-identity} forces $|\gamma(t)|$ to be constant on
$\Gamma_{\mathrm{out}}$. Hence $\Gamma_{\mathrm{out}}$ is a circle centered at the origin.

The conclusion now follows from Theorem~\ref{thm:outer-circle-implies-annulus}.
\end{proof}

\begin{remark}
This theorem is intentionally conditional. The draft \texttt{newnew.tex} tries to obtain the
analyticity of the outer boundary directly from a one-phase free-boundary argument, but that step
is not justified from its stated hypotheses. The valid takeaway is:
\begin{enumerate}[label=\textnormal{(\roman*)}]
\item if one can justify analyticity of the outer boundary by some separate argument, then the
boundary-moment asymptotics do force the outer boundary to be a circle;
\item once that outer circle is known, the rest of the connected-case conclusion is rigorous.
\end{enumerate}
\end{remark}

\section{Why the Simply Connected Proof Does Not Transfer Directly}

\begin{remark}
In Section~4 of \texttt{main.tex}, the simply connected proof defines a holomorphic function
\[
W(z)=\pi z\left(4\partial_z u(z)-\frac{\Phi(|z|^2)}{z}+\frac{c}{\pi z}\right)
\]
and uses that $W|_{\partial U}$ is real-valued. For simply connected $U$, this implies
$\Im W\equiv 0$ in $U$ by the maximum principle.

For connected multiply connected $U$, this implication is false if one only knows reality on the
outer boundary. The inner boundary components contribute additional boundary data for the harmonic
function $\Im W$, and there is no reason for those traces to vanish. This is the precise place
where one must not implicitly reuse the simply connected argument.
\end{remark}

\end{document}
